\section{Shrinking viable clouds and coherent approximate names}
\label{sec:shrinking}

\subsection{When direct stability is unavailable}

Suppose observations are controlled only along compatible edges, and different
finite paths can yield different candidates at the same level.  Finite
amalgamation produces a path for each chosen horizon, but horizons may choose
unrelated paths.  A coherent approximate name can still be extracted if all
candidates that participate in a sufficiently long full path have small
observational diameter.

Let \(d\) be a coded rational pseudometric, and suppose every compatible edge
satisfies
\[
 d(\Obs(p,k,a),\Obs(p,k+1,b))\le\delta(p,k).
\]
Choose a nondecreasing horizon \(H(p,k)\ge k\).  Define the full-viable cloud
\[
 V_{p,k}=\{a<B(p,k):\FullV(p,H(p,k),k,a)=1\}.
\]
Let \(s(p,k)\) be its least member and
\(x(p,k)=\Obs(p,k,s(p,k))\).

\begin{theorem}[Shrinking viable-cloud theorem]
\label{thm:shrinking}
Assume iPRA proves:
\begin{enumerate}[label=(SV\arabic*)]
\item every \(V_{p,k}\) is nonempty;
\item \(H(p,k)\le H(p,\ell)\) when \(k\le\ell\);
\item every two members of \(V_{p,k}\) are at distance at most
      \(\varepsilon(p,k)\); and
\item every compatible edge has motion at most \(\delta(p,k)\).
\end{enumerate}
Then for \(k\le\ell\),
\[
 d(x(p,k),x(p,\ell))
 \le \varepsilon(p,k)+
      \sum_{j=k}^{\ell-1}\delta(p,j).
\]
If an explicit tail modulus makes the right-hand side at most \(2^{-n}\), the
selected observations form a regular name.
\end{theorem}

\begin{proof}
Choose, by bounded search from the full-viability certificate, a complete path
to horizon \(H(p,\ell)\) passing through \(s(p,\ell)\).  Let \(b\) be its
level-\(k\) member.  Monotonicity of horizons permits truncation of this path at
\(H(p,k)\), so \(b\in V_{p,k}\).  The cloud-diameter hypothesis gives
\[
 d(\Obs(p,k,s(p,k)),\Obs(p,k,b))\le\varepsilon(p,k).
\]
Telescoping the edge bounds along the path segment from \(k\) to \(\ell\)
gives
\[
 d(\Obs(p,k,b),\Obs(p,\ell,s(p,\ell)))
 \le\sum_{j=k}^{\ell-1}\delta(p,j).
\]
One triangle inequality yields the result.

The object proof uses only finite objects: a least path code, path truncation,
a Boolean cloud-membership proof, and quantifier-free induction for the finite
sum of edge bounds.  The final regularity theorem substitutes the supplied tail
modulus; no infinite path is selected.
\end{proof}

\subsection{Sources of shrinking}

The cloud-diameter hypothesis is quantitative determinacy.  It can arise from:

\begin{itemize}
\item \textbf{isolating balls:} every viable candidate lies in one rational ball;
\item \textbf{strong convexity:} a residual bound and a positive convexity
      constant control distance to the unique minimizer;
\item \textbf{residual-to-distance estimates:} a certified error bound controls
      distance to an isolated zero or cluster;
\item \textbf{algebraic separation:} exact coefficient data isolate disjoint
      root clusters;
\item \textbf{contraction:} repeated certified contraction shrinks all viable
      candidates in one box; or
\item \textbf{finite Hausdorff control:} two candidate clouds approximate the
      same compact object in a finite coded metric.
\end{itemize}

Qualitative uniqueness alone supplies none of these rates.  The modulus is a
mathematical premise, not a logical consequence of uniqueness.

\subsection{Exact labels under a gap}

Suppose candidates carry labels from a finite set, and candidates with
different labels are at distance at least a positive rational \(g\).  Once the
shrinking bound is below \(g\), selected labels cannot differ.

\begin{corollary}[Eventual label stability]
\label{cor:label-gap}
If a primitive-recursive threshold \(K(p)\) proves
\[
 \varepsilon(p,k)+\sum_{j=k}^{m-1}\delta(p,j)<g
 \quad(k\ge K(p),\;m\ge k),
\]
then the selected labels are constant beyond \(K(p)\), and an exact label is
obtained by evaluating at \(K(p)\).
\end{corollary}

This is the sound interface for ordering separated roots, choosing a pivot,
selecting a chart, fixing an orientation, or identifying a branch.  At a
collision or zero gap, the theorem correctly returns only the unlabeled or
relational object.

\subsection{Comparison with the stable-observable shortcut}

The two methods solve different problems.

\begin{center}
\begin{tabularx}{0.94\textwidth}{@{}lXX@{}}
\toprule
 & Stable-observable shortcut & Shrinking viability \\
\midrule
Input comparison & Any two certificates, possibly independent & Compatible
paths plus same-level viable-cloud diameter \\
Need paths? & No & Yes, but only finite paths \\
Best use & Canonical aggregate or invariant & Refining local approximations \\
Modulus & Direct cross-precision bound & Diameter plus summable edge motion \\
Failure mode & Observable itself is unstable & Viable clouds do not shrink \\
\bottomrule
\end{tabularx}
\end{center}

An application should attempt direct stability first and introduce viability
only when the mathematics genuinely depends on refinement paths.
